\section{Limitations}
The modern seven-sequence VisDrone result is a documented second validation
read: a static control had already read those sequences before AC--MOT v3. It is
therefore a frozen comparison, not an untouched held-out confirmation; the
historical test-dev experiment is the cleaner historical held-out evidence.
The freeze also records that Trial 7 is not the literal argmax of the
documented modern selection rule under either archived eligibility
interpretation. We report this provenance limitation without rerunning the
search.

Matched-static and shuffled-timing controls weaken a claim that exact switching
timing causes the accuracy gain. The project class-agnostic evaluator is not an
official VisDrone leaderboard protocol. Only one modern detector checkpoint is
used: a third-party VisDrone-DET-trained YOLO11m model rather than a
VisDrone-MOT-trained or published OATrack detector. Two tracker hosts show
interaction but not tracker independence, and the OATrack implementation is a
controlled repository reimplementation rather than a direct published-system
comparison.

The historical V1/V2 studies use older detectors, custom held-out protocols,
and objective-specific profiles. The U2MOT transfer uses a separate VisDrone
pipeline with unsupervised boundary recalibration and no HOTA or confidence
interval; the SparseTrack transfer is a same-development-set, near-neutral
MOT17 study whose bootstrap interval includes zero. Neither establishes
universal cross-pipeline quality. UAVDT transfer is a single frozen-policy run
per host without cross-validation or bootstrap, and OATrack HOTA is essentially
flat there.

T4 timing is measured, but image-read work is not instrumented identically in
the baseline and AC--MOT v3 harnesses. Finally, throughput is below 30 FPS. These
limits constrain generalization and attribution while leaving a reproducible
detector-operating-point study with measured quality, cost, and transfer
behaviour.
